% Agent documentation:
% Purpose: define the marketing logic that the phase plan turns into concrete
% actions.
% Downstream use: frames the Nesodden and Bærum/Asker phases.
% Revision guardrails: keep this section short. Concrete channel and campaign
% choices belong in the phase sections.
\subsection{Tiltak}

Tiltakene skal gjøre posisjoneringen konkret i kundereisen. Zipline skal ikke
bare kommunisere at dronelevering er ny teknologi. Tjenesten må først
oppfattes som et trygt og enkelt alternativ til vanlig hjemlevering. Deretter
kan Zipline vise at leveringen er raskere og rimeligere enn Foodora og Wolt.

Tiltakene styres av tre KPI-er. Den første er kostnad per ny kunde (CAC). Den
andre er konverteringsrate fra kjennskap til første kjøp. Den tredje er
snittkurv per ordre. Nesodden brukes til å teste tillit, lokal forankring og
lav prøveterskel. Bærum og Asker brukes til å teste konkurranse mot etablerte
leveringsplattformer.

\subsubsection{Posisjonering i tiltakene}

Likhetspunktene må komme før differensieringspunktene. Kunden må forstå at
Zipline har enkel bestilling, relevant vareutvalg, sanntidssporing og
pålitelig levering. Dette er grunnkrav i kategorien hjemlevering. Hvis disse
kravene er uklare, vil pris og hastighet ha liten effekt fordi tjenesten
fortsatt oppleves som usikker.

Når likhetspunktene er troverdige, bør Zipline vise pris- og tidsfordelen
direkte. Budskapet bør være enkelt. Kunden får samme type bestilling som hun
allerede kjenner, men med lavere leveringsgebyr og kortere leveringstid.
Dette bør vises i annonser, partnerflater og appen. Poenget er ikke å selge
``fremtidens levering''. Poenget er å gjøre neste bestilling enklere,
raskere og billigere.

Hotellings modell og medianvelgerteoremet viser hvorfor Zipline ikke bør
gjøre differensieringen for abstrakt. Konkurrenter i modne markeder trekkes
ofte mot midten fordi de forsøker å nå flest mulig kunder
\parencite[Forelesning 5]{aspelund2026_forelesninger}. Zipline må derfor
være annerledes på de få egenskapene som faktisk betyr noe i
kjøpsøyeblikket. For denne planen er det pris og hastighet. Hvis tjenesten
fremstår for fremmed, taper Zipline på opplevd risiko. Hvis den fremstår som
samme bestillingsopplevelse med lavere pris og raskere levering, blir
differensieringen lettere å velge.

De konkrete tiltakene legges derfor inn i faseplanen. Det gjør at tiltakene
knyttes direkte til markedet de skal virke i.
